\documentclass[leqno]{article}
% Shared mark scheme preamble for resources.danieljsmith.org: the pack
% preamble, unchanged, then the scheme's own layer. A scheme's .tex is
% rendered from its YAML (scripts/README.md, Mark schemes), never edited.
\input{\CurrentFilePath/pack_preamble.tex}

% The page: equal side margins, so the table runs the full width of the
% header rule. The footers' right edge follows the text's.
\geometry{left=0.8in, right=0.8in}
\makeatletter
\renewcommand{\djs@footR}[1]{%
  \put(\dimexpr1in+\hoffset+\oddsidemargin+\textwidth\relax,
    -\dimexpr\paperheight-0.32in\relax){\makebox[0pt][r]{#1}}}
\makeatother

\usepackage[export]{adjustbox}
\usepackage{tabularray}
\definecolor{msGrid}{RGB}{176, 172, 164}
\definecolor{msTint}{RGB}{245, 243, 239}

% The header on every page, the front page included, and the booklet label
% under the contents-page title. \djsQuestionPageStyle stays empty.
\newcommand{\djsSchemeHeader}[1]{%
  \djsPackHeader[markscheme]{\textit{Mark scheme} - #1}%
  \renewcommand{\djsPackTitle}{#1}%
  \renewcommand{\djsBookletLabel}{Mark scheme}%
}

% The front page: the contents heading, spaced as tocloft spaces the packs'
% "Questions", then the contents table. A question's name and page number
% both link to its first page.
\newcommand{\msContentsHeading}{%
  \vskip\cftbeforetoctitleskip
  {\noindent\cfttoctitlefont\contentsname\par}%
  \vskip\cftaftertoctitleskip}
\NewTblrEnviron{mscontents}
\SetTblrInner[mscontents]{
  colspec = {Q[wd=26mm, l] X[l] Q[wd=16mm, c] Q[wd=14mm, r]},
  width = \linewidth,
  rowsep = 4pt, colsep = 6pt,
  hlines = {0.4pt, msGrid}, vlines = {0.4pt, msGrid},
  hline{2} = {2.2pt, scarlet},
  row{1} = {font=\small\bfseries},
  row{Z} = {font=\bfseries},
  column{2} = {font=\small},
}
\newcommand{\msContentsQuestion}[1]{\hyperref[q:#1]{Question #1}}
\newcommand{\msContentsPage}[1]{\hyperref[q:#1]{\pageref*{q:#1}}}
\newcommand{\msSep}{\hspace{0.9em}}

% The conventions pages' layout: headings in the documents' own bold, tables
% ruled and padded as the scheme's, roomy lists.
\newcommand{\msHead}[1]{%
  \par\vspace{18pt}%
  {\noindent\bfseries #1\par}%
  \vspace{6pt}}
\NewTblrEnviron{mstblr}
\SetTblrInner[mstblr]{rowsep = 6pt, colsep = 6pt,
  hlines = {0.4pt, msGrid}, vlines = {0.4pt, msGrid}}
\newenvironment{msList}
  {\begin{itemize}[leftmargin=1.2em, itemsep=5pt, topsep=3pt, parsep=0pt]}
  {\end{itemize}}

% A question's first page: the link target and the PDF bookmark, as
% \questionitem gives the packs.
\newcommand{\msQnum}{}
\newcommand{\msQuestion}[1]{%
  \clearpage
  \renewcommand{\msQnum}{#1}%
  \phantomsection\label{q:#1}%
  \addcontentsline{toc}{subsection}{\protect\djsTocEntry{#1}}}

% The scheme's table: Question 16mm, Scheme the rest, Marks 16mm, Notes
% 44mm; every cell ruled in msGrid; the scarlet rule under the head row. A
% long table has no caption or "continued" foot; a continuation page opens
% with a muted line above the repeated head row. The rules across a part's
% rows, the alternative's tint and the total row are the rendered table's.
\NewTblrTheme{ms}{%
  \DefTblrTemplate{firsthead}{default}{}%
  \DefTblrTemplate{middlehead,lasthead}{default}{%
    {\small\itshape\color{djsMuted}Question \msQnum\ continued\par}%
    \vspace{6pt}}%
  \DefTblrTemplate{firstfoot,middlefoot,lastfoot}{default}{}%
}
\NewTblrEnviron{msscheme}
\SetTblrOuter[msscheme]{long, theme=ms, entry=none, label=none}
\SetTblrInner[msscheme]{
  colspec = {Q[wd=16mm, l] X[l] Q[wd=16mm, c] Q[wd=44mm, l]},
  width = \linewidth,
  rowsep = 6pt, colsep = 6pt,
  rowhead = 1,
  vlines = {0.4pt, msGrid},
  hline{1} = {0.4pt, msGrid}, hline{2} = {2.2pt, scarlet},
  hline{Z} = {0.4pt, msGrid},
  column{1} = {font=\bfseries}, column{3} = {font=\bfseries},
  column{4} = {font=\small},
  row{1} = {rowsep=4pt, font=\small\bfseries},
}

% Two marks on a row: each note starts with its code.
\newcommand{\msNote}[2]{\textbf{#1}\enspace #2}
% A full-row figure: the worked solution's own picture, scaled uniformly and
% only down, to the table less the Question column. The picture is drawn
% with the line width it has in the solutions document, at its list depth
% there (the optional argument, 1 by default: a part's working), so one
% sized by \linewidth keeps its size; tabularray sets the cell's \linewidth
% to the cell's measured width, which would shrink it at every pass. The
% widths are measured once, in a solution box at the packs' text width
% (pack_preamble.tex's margins, 0.8in and 1in).
\newlength{\msSpanW}
\makeatletter
\newcommand{\ms@keepLW}[1]{%
  \expandafter\xdef\csname ms@solnLW@#1\endcsname{\the\linewidth}}
\AtBeginDocument{%
  \setlength{\msSpanW}{\dimexpr\textwidth-16mm-28pt\relax}%
  \setbox\z@\vbox{%
    \begin{minipage}{\dimexpr\paperwidth-1.8in\relax}
      \begin{djsSolution}
        \ms@keepLW{0}%
        \begin{enumerate}\item \ms@keepLW{1}%
        \begin{enumerate}\item \ms@keepLW{2}%
        \begin{enumerate}\item \ms@keepLW{3}%
        \end{enumerate}\end{enumerate}\end{enumerate}
      \end{djsSolution}
    \end{minipage}}}
\newcommand{\msFigure}[2][1]{%
  \adjustbox{max width=\msSpanW, max totalheight=0.6\textheight, valign=t}{%
    \hsize=\csname ms@solnLW@#1\endcsname
    \linewidth=\hsize
    #2}}
\makeatother
% A figure mark's feature, in the Scheme column.
\newcommand{\msFeature}[1]{\textit{#1}}
% An alternative's title row.
\newcommand{\msAlternative}[1]{\textit{Alternative:} #1}
% Drafts only: a queried mark's code, the query's why under its note, and a
% flag in its part's place.
\newcommand{\msQueried}[2]{#1~\(\to\) #2?}
\newcommand{\msQuery}[1]{{\itshape\color{djsMuted}Query: #1}}
\newcommand{\msFlag}[2]{\textbf{#1:}\enspace #2}


\djsSchemeHeader{Shortest Distances}
\djsSitePath{Core Pure}{Vectors, Matrices \& Linear Transformations}

\djsLastUpdated{25 October 2026}
\begin{document}
\djsSetQuestionCount{8}
\djsAddMarks{1}{13}
\djsAddMarks{2}{18}
\djsAddMarks{3}{8}
\djsAddMarks{4}{10}
\djsAddMarks{5}{8}
\djsAddMarks{6}{17}
\djsAddMarks{7}{9}
\djsAddMarks{8}{16}
\vspace*{6pt}
\djsContentsTitle
\vspace{30pt}
\msContentsHeading
\begin{mscontents}{}
Question & Marks by part & Marks & Page \\
\msContentsQuestion{1} & (a)\,2\msSep (b)\,2\msSep (c)\,3\msSep (d)\,3\msSep (e)\,3 & 13 & \msContentsPage{1} \\
\msContentsQuestion{2} & (a)\,4\msSep (b)\,5\msSep (c)(i)\,2\msSep (c)(ii)\,3\msSep (c)(iii)\,4 & 18 & \msContentsPage{2} \\
\msContentsQuestion{3} & (a)\,2\msSep (b)\,2\msSep (c)\,4 & 8 & \msContentsPage{3} \\
\msContentsQuestion{4} & (a)\,6\msSep (b)\,4 & 10 & \msContentsPage{4} \\
\msContentsQuestion{5} & (i)\,6\msSep (ii)\,2 & 8 & \msContentsPage{5} \\
\msContentsQuestion{6} & (a)\,3\msSep (b)\,3\msSep (c)\,2\msSep (d)\,4\msSep (e)(i)\,3\msSep (e)(ii)\,2 & 17 & \msContentsPage{6} \\
\msContentsQuestion{7} & (a)(i)\,1\msSep (a)(ii)\,2\msSep (b)(i)\,3\msSep (b)(ii)\,3 & 9 & \msContentsPage{7} \\
\msContentsQuestion{8} & (a)\,7\msSep (b)\,4\msSep (c)\,5 & 16 & \msContentsPage{8} \\
\SetCell[c=2]{r} Total & & 99 & \\
\end{mscontents}
\newpage
% The conventions pages, shared by every pack's mark scheme: one source of
% truth, \input after the front page, with no strand or pack logic. The keys
% on the first side, the rules on the second. The scheme layer defines
% \msHead, msList and mstblr.
\vspace*{6pt}%
{\noindent\cfttoctitlefont Reading the mark scheme\par}%
\vspace{10pt}%
These rules apply throughout, unless a note says otherwise.

\msHead{Marks}
Every mark is worth one, so \textbf{M1 A1} on a row is two marks. Only marks
coded ft follow through.\par
\vspace{10pt}%
\begin{mstblr}{
  colspec = {Q[wd=15mm, l] X[l]},
  width = \linewidth,
  hline{2} = {2.2pt, scarlet},
  column{1} = {font=\bfseries},
  row{1} = {font=\small\bfseries, rowsep=4pt},
}
Code & Meaning \\
M1 & Method: a correct method, applied to this question. A slip doesn't cost it if the method is clear. A formula must be used, not just quoted. \\
dM1 & Dependent method: as M1, but only if the method mark before it was earned, unless its note names another. \\
A1 & Accuracy: a correct answer or step. It needs the method mark it follows. \\
A1ft & As A1, following through from an earlier answer that the note names. \\
A1~AG & The last mark of a part whose answer is printed in the question: complete working, with no errors. \\
B1 & A correct value, fact or statement. It needs no method mark. \\
B1ft & As B1, following through from an earlier answer that the note names. \\
E1 & A written explanation, interpretation or conclusion. Concluding a proof by induction needs every earlier mark. \\
\end{mstblr}

\msHead{Abbreviations}
ft and AG qualify a code in the Marks column. The others appear in the
notes.\par
\vspace{10pt}%
\begin{mstblr}{
  colspec = {Q[wd=17mm, l] X[l] Q[wd=15mm, l] X[l]},
  width = \linewidth,
  column{1} = {font=\bfseries}, column{3} = {font=\bfseries},
}
ft & follow through & AG & answer given in the question \\
cao & correct answer only, with no follow-through & cso & correct solution only: no errors in the part \\
oe & or an equivalent form & awrt & any answer that rounds to this value \\
soi & seen, or implied by later working & isw & ignore working after a correct answer \\
SC & special case: credit outside the main method & their & the student's own earlier value \\
condone & a slip that doesn't cost the mark & [\;\,] & working the mark doesn't need \\
s.f. & significant figures & d.p. & decimal places \\
\end{mstblr}

\newpage
\msHead{Reading the scheme}
\begin{msList}
\item Each row is a line of working, with the marks it earns beside it.
  Where a row carries two marks, each note starts with its code.
\item \textbf{[\textit{n}]} ends each part with its total.
\item A shaded \textit{Alternative} is a second method for the same total.
  Mark an answer by the method it follows.
\item Any other valid method earns the equivalent marks, unless the question
  names a method or says \textit{hence}.
\item A sketch earns one B1 for each feature named below its figure.
\end{msList}

\msHead{Accuracy and answers}
\begin{msList}
\item Unless a question says otherwise, give answers to 3 significant
  figures, and angles in degrees to 1 decimal place. Angles in radians take
  3 significant figures.
\item A final answer with more figures is accepted if it rounds to the
  scheme's value; one with fewer is not. An answer spoiled by rounding too
  early loses its accuracy mark.
\item Where a question asks for an exact answer, a decimal loses the final
  accuracy mark. Otherwise either is accepted.
\item Wrong or missing units cost nothing unless the question asks for them.
\item A correct answer with no working earns full marks, unless the question
  needs the working: show that, prove, hence, a method it names, or an exact
  answer. A wrong answer with no working earns nothing. Calculator answers
  are treated the same way.
\item A correct answer, once reached, isn't spoiled by later working (isw).
\end{msList}

\msHead{Value of $\boldsymbol{g}$}
\begin{msList}
\item Take $g = 9.8\,\mathrm{m\,s^{-2}}$ unless a question states a value,
  which is then the only one accepted. Answers that use $g = 9.8$ may be
  given to 2 or 3 significant figures.
\item Where a question doesn't state $g$, an answer worked with $g = 10$
  throughout earns the same marks, keeping any value printed in the
  question. The worked solutions list the final answers with $g = 10$.
\end{msList}


\msQuestion{1}
\begin{msscheme}{
  hline{3-4,6-7,9-11,13-15,17-19} = {2-Z}{0.4pt, msGrid},
  hline{5,8,12,16,20} = {0.4pt, msGrid},
  row{20} = {font=\bfseries},
}
Question & Scheme & Marks & Notes \\
1(a) & {\(\frac{\lvert 3(5) + 1(-5) + 2(3) + 4\rvert}{\sqrt{3^2 + 1^2 + 2^2}}\)} & M1 & {Substitutes \(P\) and the plane into the distance formula, or finds the foot of the perpendicular. Condone one sign slip.} \\
& {\(= \frac{20}{\sqrt{14}}\ \left[= \frac{10\sqrt{14}}{7}\right]\)} & A1 & {\(\frac{20}{\sqrt{14}}\) oe, or awrt 5.35.} \\
\nopagebreak & & [2] & \\
1(b) & {\(2 + \lambda = 5 \Rightarrow \lambda = 3\)} & M1 & {Equates one component of the general point of \(L\) to that of \(P\) and solves for \(\lambda\).} \\
& {\(y = 1 - 2(3) = -5\), \(z = -3 + 2(3) = 3\)\\[4pt] so \(P\) lies on \(L\)} & A1~AG & {Both other components shown to match \(P\) at \(\lambda = 3\), or \(\lambda = 3\) from all three, with a conclusion.} \\
\nopagebreak & & [2] & \\
1(c) & {\(3(2 + \lambda) + (1 - 2\lambda) + 2(-3 + 2\lambda) = -4\)\\[4pt] \(\Rightarrow \lambda = \ldots\)} & M1 & {Substitutes all three components of \(L\) into the plane equation and solves for \(\lambda\). Condone sign slips.} \\
& {\(\lambda = \their{{-1}} \Rightarrow (2 - 1,\ 1 + 2,\ -3 - 2)\)} & dM1 & {Substitutes their \(\lambda\) into the equation of \(L\).} \\
& {\((1, 3, -5)\)} & A1 & \\
\nopagebreak & & [3] & \\
1(d) & {\(\mathbf{u}\cdot\mathbf{n} = 3 - 2 + 4 = 5\), \(\lvert\mathbf{u}\rvert = 3\), \(\lvert\mathbf{n}\rvert = \sqrt{14}\)} & B1 & {All three values, seen or used.} \\
& {\(\sin\theta = \frac{\their{5}}{\their{3}\,\their{\sqrt{14}}} \Rightarrow \theta = \ldots\)} & M1 & {Uses \(\sin\theta = \frac{\lvert\mathbf{u}\cdot\mathbf{n}\rvert}{\lvert\mathbf{u}\rvert\lvert\mathbf{n}\rvert}\) with their values, or finds the angle with the normal by cosine and subtracts it from \(90^\circ\).} \\
& {\(\theta = 26.5^\circ\)} & A1 & {awrt \(26.5^\circ\), or awrt 0.462 radians.} \\
\nopagebreak & & [3] & \\
1(e) & {\(PC = \left(3 - (-1)\right) \times 3 = \ldots\ [= 12]\)} & M1 & {Finds \(PC\) as the change in \(\lambda\) times their \(\lvert\mathbf{u}\rvert\), or as the magnitude of \(\overrightarrow{PC}\) using their point from (c).} \\
& {\(\their{12} \times \their{\frac{5}{3\sqrt{14}}}\)} & dM1 & {Multiplies their \(PC\) by their \(\sin\theta\) from (d), or by the cosine of their angle with the normal.} \\
& {\(= \frac{20}{\sqrt{14}}\), as in (a)} & A1 & {\(\frac{20}{\sqrt{14}}\) oe, or awrt 5.35, with a statement that it agrees with (a).} \\
\nopagebreak & & [3] & \\
\nopagebreak \SetCell[c=2]{r} Total for Question 1 & & 13 & \\
\end{msscheme}

\msQuestion{2}
\begin{msscheme}{
  hline{3-6,8-12,14-15,17-19,21-24} = {2-Z}{0.4pt, msGrid},
  hline{7,13,16,20,25} = {0.4pt, msGrid},
  row{25} = {font=\bfseries},
}
Question & Scheme & Marks & Notes \\
2(a) & {\(\overrightarrow{CD} = \begin{pmatrix}3\\3\\-3\end{pmatrix}\)} & B1 & {Or any nonzero multiple, such as \(\mathbf i + \mathbf j - \mathbf k\).} \\
& {\(\sin\alpha = \frac{\lvert(1)(1)+(1)(1)+(-1)(1)\rvert}{\sqrt3\sqrt3}\)} & M1 & {Uses the scalar product of their direction of \(CD\) and \(\mathbf i + \mathbf j + \mathbf k\) over the product of magnitudes, as \(\sin\alpha\), or as a cosine then subtracts from \(90^\circ\).} \\
& {\(\sin\alpha = \frac13\)} & A1 & {Or the cosine of the angle with the normal \(= \frac13\).} \\
& {\(\alpha = 19.5^\circ\)} & A1 & \\
\nopagebreak & & [4] & \\
2(b) & {\(\mathbf r = \begin{pmatrix}0\\0\\2\end{pmatrix} + t\begin{pmatrix}1\\1\\-1\end{pmatrix}\)} & M1 & {Uses \(C\) or \(D\) and their direction of \(CD\) from (a) in a vector or parametric equation.} \\
& {\((x, y, z) = (t, t, 2 - t)\)} & A1 & {oe for their point and parameter.} \\
& {\(P_1\colon\ t + t + 2 - t = 3 \Rightarrow t = \ldots\)\\[4pt] \(P_2\colon\ 2t - t = 1 \Rightarrow t = \ldots\)} & M1 & {Substitutes their general point into both plane equations and solves each, or into one and tests the resulting point in the other.} \\
& {\(t = 1\) from both, so the point is common to \(P_1\) and \(P_2\)} & A1 & {The same parameter from both, or the point from one shown to satisfy the other, with a conclusion that it is common. A0 if the conclusion is missing.} \\
& {\((1, 1, 1)\)} & A1 & \\
\nopagebreak & & [5] & \\
2(c)(i) & {\((\mathbf i + \mathbf j + \mathbf k) \times (2\mathbf i - \mathbf j)\)\\[4pt] \(\phantom{(} = \mathbf i(0 + 1) - \mathbf j(0 - 2) + \mathbf k(-1 - 2)\)} & M1 & {Expands a determinant or uses the component formula for all three components. Condone one sign slip.} \\
& {\(= \mathbf i + 2\mathbf j - 3\mathbf k\)} & A1 & \\
\nopagebreak & & [2] & \\
2(c)(ii) & {\(\lvert\mathbf i + 2\mathbf j - 3\mathbf k\rvert = \sqrt{14}\)\\[4pt] \(\lvert\mathbf i + \mathbf j + \mathbf k\rvert = \sqrt3,\ \lvert 2\mathbf i - \mathbf j\rvert = \sqrt5\)} & B1 & {All three needed.} \\
& {\(\sin\theta = \frac{\their{\sqrt{14}}}{\sqrt3\sqrt5} \Rightarrow \theta = \ldots\)} & M1 & {Divides the magnitude of their vector from (c)(i) by the product of the normals' magnitudes as \(\sin\theta\), and finds \(\theta\). Not for a scalar-product route that doesn't use (c)(i).} \\
& {\(\theta = 75.0^\circ\)} & A1 & \\
\nopagebreak & & [3] & \\
2(c)(iii) & {\(\begin{pmatrix}0\\0\\2\end{pmatrix} - \begin{pmatrix}1\\1\\1\end{pmatrix} = \begin{pmatrix}-1\\-1\\1\end{pmatrix}\)\\[4pt] \(\begin{pmatrix}-1\\-1\\1\end{pmatrix} \times \begin{pmatrix}1\\2\\-3\end{pmatrix} = \ldots\)} & M1 & {Forms the vector from a point on both planes, such as their point from (b), to \(C\), and crosses it with their direction from (c)(i). Condone sign slips.} \\
& {\(\begin{pmatrix}1\\-2\\-1\end{pmatrix}\)} & A1 & {Or its negative.} \\
& {\(\text{dist} = \frac{\their{\sqrt6}}{\sqrt{14}}\)} & dM1 & {Divides the magnitude of their cross product by the magnitude of their direction.} \\
& {\(= \sqrt{\tfrac37} \approx 0.655\)} & A1 & \\
\nopagebreak & & [4] & \\
\nopagebreak \SetCell[c=2]{r} Total for Question 2 & & 18 & \\
\end{msscheme}

\msQuestion{3}
\begin{msscheme}{
  hline{3-4,6-7,9-12} = {2-Z}{0.4pt, msGrid},
  hline{5,8,13} = {0.4pt, msGrid},
  row{13} = {font=\bfseries},
}
Question & Scheme & Marks & Notes \\
3(a) & {\((1)(1) + (-2)(2) + (3)(1) = \ldots\)\\[4pt] \((1)(4) + (-2)(-1) + (3)(-2) = \ldots\)} & M1 & {Forms the scalar product of the vector with both direction vectors, or finds the vector product of the direction vectors.} \\
& {\(1 - 4 + 3 = 0\) and \(4 + 2 - 6 = 0\)\\[4pt] Perpendicular to both direction vectors \(\Rightarrow\) perpendicular to \(\Pi\)} & A1~AG & {Both products shown evaluated to zero, or the vector product shown to be a multiple of the vector, with a conclusion.} \\
\nopagebreak & & [2] & \\
3(b) & {\(x - 2y + 3z = (1)(2) + (-2)(-1) + (3)(3) = \ldots\)} & M1 & {Substitutes \((2, -1, 3)\), or another point of \(\Pi\), into \(x - 2y + 3z = d\), or expands \((\mathbf{i} - 2\mathbf{j} + 3\mathbf{k})\cdot(\mathbf{r} - \mathbf{a}) = 0\). Condone sign slips.} \\
& {\(x - 2y + 3z = 13\)} & A1 & {Cartesian form, oe. \(\mathbf{r}\cdot(\mathbf{i} - 2\mathbf{j} + 3\mathbf{k}) = 13\) alone scores A0.} \\
\nopagebreak & & [2] & \\
3(c) & {\(A = (-2 + t,\ 1 + t,\ 4 + t)\)\\[4pt] \(\frac{|(-2 + t) - 2(1 + t) + 3(4 + t) - \their{13}|}{\sqrt{1 + 4 + 9}}\)\\[4pt] \(\phantom{A} = \frac{9}{\sqrt{14}}\)} & M1 & {Uses a general point of \(\ell\) in a point-to-plane distance with their plane from (b), set equal to \(\frac{9}{\sqrt{14}}\), or an equivalent method such as parallel planes. Condone sign slips.} \\
& {\(|2t - 5| = 9\)} & A1 & {oe, such as \((2t - 5)^2 = 81\) or \(2t - 5 = \pm 9\).} \\
& {\(2t - 5 = 9\) or \(2t - 5 = -9 \Rightarrow t = \ldots\ [= 7, -2]\)} & dM1 & {Solves both cases, or their quadratic, for two values of \(t\).} \\
& {\(A = (5, 8, 11)\) or \((-4, -1, 2)\)} & A1 & {Both points, as coordinates or position vectors, and no others.} \\
\nopagebreak & & [4] & \\
\nopagebreak \SetCell[c=2]{r} Total for Question 3 & & 8 & \\
\end{msscheme}

\msQuestion{4}
\begin{msscheme}{
  hline{3-8,10-13} = {2-Z}{0.4pt, msGrid},
  hline{9,14} = {0.4pt, msGrid},
  row{14} = {font=\bfseries},
}
Question & Scheme & Marks & Notes \\
4(a) & {\(\overrightarrow{AB} = \begin{pmatrix}1\\2\\0\end{pmatrix}\), \(\overrightarrow{AC} = \begin{pmatrix}1\\0\\-1\end{pmatrix}\)\\[4pt] \(\overrightarrow{AB}\times\overrightarrow{AC} = \ldots\)} & M1 & {Finds two non-parallel vectors in the plane and attempts their cross product, or solves two scalar products with them equal to zero. Condone sign slips.} \\
& {\(\mathbf n = \begin{pmatrix}-2\\1\\-2\end{pmatrix}\)} & A1 & {Any nonzero multiple.} \\
& {\(2x - y + 2z = 2(1) - (-2) + 2(3)\)} & M1 & {Substitutes a point of \(\Pi\) into \(ax + by + cz = d\) or \(\mathbf r\cdot\mathbf n = d\), with their normal, to find \(d\).} \\
& {\(2x - y + 2z = 10\)} & A1 & {oe, including \(\mathbf r\cdot\begin{pmatrix}2\\-1\\2\end{pmatrix} = 10\) or any multiple.} \\
& {\(\text{dist} = \frac{\lvert 2(4) - 1 + 2(0) - 10\rvert}{\sqrt{2^2 + (-1)^2 + 2^2}} = \ldots\)} & M1 & {Uses the point-to-plane formula with \(D\) and their plane \(ax + by + cz = d\), or projects \(\overrightarrow{AD}\) onto their normal. Condone sign slips.} \\
& {\(= 1\)} & A1 & \\
\nopagebreak & & [6] & \\
4(b) & {\(\begin{pmatrix}2\\1\\2\end{pmatrix}\times\begin{pmatrix}1\\-2\\1\end{pmatrix} = \ldots\)} & M1 & {Attempts the cross product of the two direction vectors, or solves two scalar products with them equal to zero. Condone sign slips.} \\
& {\(= \begin{pmatrix}5\\0\\-5\end{pmatrix}\)} & A1 & {Any nonzero multiple.} \\
& {\(\begin{pmatrix}2\\-3\\5\end{pmatrix}\cdot\their{\begin{pmatrix}5\\0\\-5\end{pmatrix}}\ [= -15]\)\\[4pt] \(\text{dist} = \frac{15}{\sqrt{5^2 + 0^2 + (-5)^2}} = \ldots\)} & dM1 & {Divides the scalar product of a vector joining the lines with their normal by the normal's magnitude. Condone a missing modulus.} \\
& {\(= \frac{3\sqrt2}{2}\)} & A1 & {Any exact form, such as \(\frac{3}{\sqrt2}\).} \\
\nopagebreak & & [4] & \\
\nopagebreak \SetCell[c=2]{r} Total for Question 4 & & 10 & \\
\end{msscheme}

\msQuestion{5}
\begin{msscheme}{
  hline{3-8,10-11} = {2-Z}{0.4pt, msGrid},
  hline{9,12} = {0.4pt, msGrid},
  row{12} = {font=\bfseries},
}
Question & Scheme & Marks & Notes \\
5(i) & {Sets \(z = t\): \(x - y + t = 6,\ 2x + y - t = 3\)\\[4pt] \(\phantom{A} \Rightarrow x = \ldots,\ y = \ldots\)} & M1 & {Solves the plane equations with one variable as parameter, or crosses the normals and finds a common point. Condone slips.} \\
& {\(\mathbf r = \begin{pmatrix}3\\-3\\0\end{pmatrix} + t\begin{pmatrix}0\\1\\1\end{pmatrix}\)} & A1 & {Any point on \(\ell_2\) and any multiple of the direction.} \\
& {\(\begin{pmatrix}1\\2\\-1\end{pmatrix}\times\their{\begin{pmatrix}0\\1\\1\end{pmatrix}} = \ldots\)} & M1 & {Forms the vector product of the two directions, or finds a vector perpendicular to both. At least two components correct for their vectors.} \\
& {\(\phantom{A} = \begin{pmatrix}3\\-1\\1\end{pmatrix}\)} & A1 & {Or any nonzero multiple.} \\
& {\(\dfrac{\left|\begin{pmatrix}1\\-2\\-3\end{pmatrix}\cdot\their{\begin{pmatrix}3\\-1\\1\end{pmatrix}}\right|}{\sqrt{11}} = \ldots\)} & dM1 & {Divides the scalar product of a vector joining the two lines with their normal by the normal's magnitude. Condone a missing modulus.} \\
& {\(\dfrac{2}{\sqrt{11}}\)} & A1 & {oe, or awrt 0.603.} \\
\nopagebreak & & [6] & \\
5(ii) & {\(\their{\begin{pmatrix}3\\-1\\1\end{pmatrix}}\cdot\begin{pmatrix}2\\-1\\3\end{pmatrix} = \ldots\)} & M1 & {Forms \(ax + by + cz = d\) with their normal from (i) and finds \(d\) from a point on \(\ell_1\).} \\
& {\(3x - y + z = 10\)} & A1 & {Any Cartesian equivalent. Not for a vector equation alone.} \\
\nopagebreak & & [2] & \\
\nopagebreak \SetCell[c=2]{r} Total for Question 5 & & 8 & \\
\end{msscheme}

\msQuestion{6}
\begin{msscheme}{
  hline{3-5,7-9,11-12,14-17,19-21,23-24} = {2-Z}{0.4pt, msGrid},
  hline{6,10,13,18,22,25} = {0.4pt, msGrid},
  row{25} = {font=\bfseries},
}
Question & Scheme & Marks & Notes \\
6(a) & {\(\det\begin{pmatrix}1&\lambda&1\\2&-1&2\\2&1&-1\end{pmatrix}\)\\[4pt] \(\phantom{\det} = 1(1-2) - \lambda(-2-4) + 1(2+2)\)} & M1 & {Expands the determinant of the coefficient matrix by any row or column, giving an expression in \(\lambda\). Condone sign slips.} \\
& {\(\phantom{\det} = 6\lambda + 3\)} & A1 & {oe, such as \(3(2\lambda+1)\).} \\
& {\(6\lambda + 3 = 0 \Rightarrow \lambda = -\tfrac12\)\\[4pt] Unique point for all \(\lambda \neq -\tfrac12\)} & A1 & {Needs the conclusion that the point is unique for every other \(\lambda\). The case \(\lambda = -\frac12\) need not be examined further.} \\
\nopagebreak & & [3] & \\
6(b) & {\(\begin{pmatrix}x\\y\\z\end{pmatrix} = \begin{pmatrix}1&2&1\\2&-1&2\\2&1&-1\end{pmatrix}^{-1}\begin{pmatrix}4\\3\\5\end{pmatrix}\)} & M1 & {Writes the system as \(\mathbf{A}\mathbf{x} = \mathbf{b}\) with \(\lambda = 2\) and uses \(\mathbf{x} = \mathbf{A}^{-1}\mathbf{b}\). Not for solving the equations without matrices.} \\
& {\(\mathbf{A}^{-1} = \frac{1}{15}\begin{pmatrix}-1&3&5\\6&-3&0\\4&3&-5\end{pmatrix}\)} & B1 & {Any equivalent form; an inverse from a calculator is accepted.} \\
& {\(\begin{pmatrix}x\\y\\z\end{pmatrix} = \begin{pmatrix}2\\1\\0\end{pmatrix}\), \(P = (2, 1, 0)\)} & A1 & \\
\nopagebreak & & [3] & \\
6(c) & {\(\frac{x+1}{1} = \frac{y-1}{-1} = \frac{z}{1} = t\)\\[4pt] \(\Rightarrow x = t - 1,\ y = 1 - t,\ z = t\)} & M1 & {Sets the ratios equal to a parameter and writes \(x\), \(y\), \(z\) in terms of it, or reads a point and a direction from the Cartesian form. Condone sign slips.} \\
& {\(\mathbf{r} = \begin{pmatrix}-1\\1\\0\end{pmatrix} + t\begin{pmatrix}1\\-1\\1\end{pmatrix}\)} & A1 & {Any point of \(\ell\) and any multiple of its direction, with \(\mathbf{r} =\) or equivalent.} \\
\nopagebreak & & [2] & \\
6(d) & {\(\their{\begin{pmatrix}2\\1\\0\end{pmatrix}} - \begin{pmatrix}-1\\1\\0\end{pmatrix} = \begin{pmatrix}3\\0\\0\end{pmatrix}\)\\[4pt] \(\begin{pmatrix}3\\0\\0\end{pmatrix} \times \begin{pmatrix}1\\-1\\1\end{pmatrix} = \ldots\)} & M1 & {Forms the vector from any point of \(\ell\) to their \(P\) and its vector product with the direction of \(\ell\). Condone sign slips.} \\
& {\(\phantom{\begin{pmatrix}3\\0\\0\end{pmatrix}} = \begin{pmatrix}0\\-3\\-3\end{pmatrix}\)} & A1 & {oe, such as its negative, or the multiple that their multiple of the direction gives.} \\
& {\(\dfrac{\sqrt{0^2 + (-3)^2 + (-3)^2}}{\sqrt{1^2 + (-1)^2 + 1^2}} = \ldots\)} & dM1 & {Divides the magnitude of their vector product by the magnitude of their direction of \(\ell\).} \\
& {\(= \sqrt6\)} & A1 & {\(\sqrt6\) or awrt 2.45.} \\
\nopagebreak & & [4] & \\
6(e)(i) & {\(\begin{pmatrix}1\\-1\\1\end{pmatrix} \cdot \begin{pmatrix}2\\1\\-1\end{pmatrix} = 2 - 1 - 1\)} & M1 & {Forms the scalar product of the direction of \(\ell\) with the normal of the plane.} \\
& {\(\phantom{\begin{pmatrix}1\\-1\\1\end{pmatrix}} = 0\)} & A1 & \\
& {\(2(-1) + 1 - 0 = -1 \neq 5\)\\[4pt] so \(\ell\) is not in the plane, and is parallel to it} & A1~AG & {Needs the zero scalar product and a point of \(\ell\) shown not to satisfy the plane's equation, then a conclusion of parallel in words.} \\
\nopagebreak & & [3] & \\
6(e)(ii) & {\(\dfrac{\left|2(-1) + 1(1) - 1(0) - 5\right|}{\sqrt{2^2 + 1^2 + (-1)^2}} = \ldots\)} & M1 & {Substitutes any point of \(\ell\) into the point-to-plane distance formula, or uses the parallel plane through \(\ell\). Unquoted, implied only by a correct substitution.} \\
& {\(= \frac{6}{\sqrt6} = \sqrt6\)} & A1 & {\(\sqrt6\) or awrt 2.45.} \\
\nopagebreak & & [2] & \\
\nopagebreak \SetCell[c=2]{r} Total for Question 6 & & 17 & \\
\end{msscheme}

\msQuestion{7}
\begin{msscheme}{
  hline{3,5-6,8-10,12-14} = {2-Z}{0.4pt, msGrid},
  hline{4,7,11,15} = {0.4pt, msGrid},
  row{15} = {font=\bfseries},
}
Question & Scheme & Marks & Notes \\
7(a)(i) & {\(\mathbf{u}\cdot\mathbf{v} = \lambda - 4\)} & B1 & \\
\nopagebreak & & [1] & \\
7(a)(ii) & {\(\mathbf{u}\times\mathbf{v} = (2-2)\mathbf{i} - (-\lambda-2)\mathbf{j}\)\\[4pt] \(\phantom{\mathbf{u}\times\mathbf{v} =} + (\lambda+2)\mathbf{k}\)} & M1 & {Attempts all three components by a determinant or the component formula. Condone sign slips.} \\
& {\(\phantom{\mathbf{u}\times\mathbf{v}} = (\lambda+2)(\mathbf{j}+\mathbf{k})\)} & A1 & {oe, such as a column vector; need not be factorised.} \\
\nopagebreak & & [2] & \\
7(b)(i) & {\(\lvert 2\mathbf{i}-2\mathbf{j}+2\mathbf{k}\rvert = 2\sqrt3\), \(\lvert\mathbf{i}+\mathbf{j}-\mathbf{k}\rvert = \sqrt3\)} & B1 & \\
& {\(\cos\theta = \dfrac{\lvert\their{2-4}\rvert}{2\sqrt3\times\sqrt3} = \ldots\)} & M1 & {Divides their (a)(i) at \(\lambda = 2\) by the product of their magnitudes in the cosine formula. Condone a missing modulus.} \\
& {\(\theta = 70.5^\circ\)} & A1 & {awrt \(70.5^\circ\) or 1.23 rad, or \(\cos^{-1}\frac13\).} \\
\nopagebreak & & [3] & \\
7(b)(ii) & {\((\mathbf{i}-2\mathbf{j}+2\mathbf{k})\times(\mathbf{i}+\mathbf{j}-\mathbf{k}) = 3(\mathbf{j}+\mathbf{k})\)} & B1ft & {ft their (a)(ii) at \(\lambda = 1\); any nonzero multiple accepted.} \\
& {\(\dfrac{\lvert(\mathbf{i}+3\mathbf{j}+4\mathbf{k})\cdot\their{(3\mathbf{j}+3\mathbf{k})}\rvert}{\lvert\their{3\mathbf{j}+3\mathbf{k}}\rvert}\)\\[4pt] \(\phantom{A} = \ldots\)} & M1 & {Divides the scalar product of a vector joining the two lines with their common normal by its magnitude. Condone sign slips in the joining vector.} \\
& {\(\dfrac{7\sqrt2}{2}\)} & A1 & {\(\frac{7\sqrt2}{2}\) or \(\frac72\sqrt2\). Not for \(\frac{7}{\sqrt2}\) or a decimal alone.} \\
\nopagebreak & & [3] & \\
\nopagebreak \SetCell[c=2]{r} Total for Question 7 & & 9 & \\
\end{msscheme}

\msQuestion{8}
\begin{msscheme}{
  hline{3-8,10-12,14-17} = {2-Z}{0.4pt, msGrid},
  hline{9,13,18} = {0.4pt, msGrid},
  row{18} = {font=\bfseries},
}
Question & Scheme & Marks & Notes \\
8(a) & {\(\overrightarrow{AD} = \begin{pmatrix}1\\1\\m\end{pmatrix}\), \(\overrightarrow{BC} = \begin{pmatrix}-1\\1\\2\end{pmatrix}\), \(\overrightarrow{AB} = \begin{pmatrix}1\\0\\0\end{pmatrix}\)} & B1 & {All three, oe; any vector joining a point of one line to a point of the other in place of \(\overrightarrow{AB}\).} \\
& {\(\begin{pmatrix}1\\1\\m\end{pmatrix}\times\begin{pmatrix}-1\\1\\2\end{pmatrix} = \begin{pmatrix}2-m\\-2-m\\2\end{pmatrix}\)} & M1 A1 & {\msNote{M1}{Attempts the vector product of their two direction vectors. Condone sign slips.}\\ \msNote{A1}{oe, such as its negative.}} \\
& {\(\dfrac{\left|\overrightarrow{AB}\cdot\their{\mathbf{n}}\right|}{\left|\their{\mathbf{n}}\right|} = \dfrac{1}{\sqrt{14}}\), \(\mathbf{n}\) their vector product} & M1 & {Uses the scalar product of their joining vector with their vector product, divided by its magnitude, and sets it equal to \(\frac{1}{\sqrt{14}}\).} \\
& {\(\dfrac{|2-m|}{\sqrt{2m^2+12}} = \dfrac{1}{\sqrt{14}}\)} & A1 & {Correct equation oe, need not be simplified. Condone a missing modulus.} \\
& {\(14(2-m)^2 = 2m^2+12\)\\[4pt] \(\Rightarrow 3m^2-14m+11 = 0 \Rightarrow m = \ldots\)} & dM1 & {Squares and solves their three-term quadratic in \(m\) by any means, a calculator included.} \\
& {\(m = 1\) or \(m = \tfrac{11}{3}\); \(m\) is an integer, so \(m = 1\)} & A1~AG & {cso. Both roots from a correct quadratic, with \(\frac{11}{3}\) rejected as not an integer. A0 if \(m = 1\) is only verified.} \\
\nopagebreak & & [7] & \\
8(b) & {\(\overrightarrow{AC}\times\overrightarrow{AD} = \begin{pmatrix}0\\1\\2\end{pmatrix}\times\begin{pmatrix}1\\1\\1\end{pmatrix} = \begin{pmatrix}-1\\2\\-1\end{pmatrix}\)} & M1 A1 & {\msNote{M1}{Attempts the vector product of a vector from a point on the line to \(C\) with the line's direction, using \(m = 1\). Condone sign slips.}\\ \msNote{A1}{Or its negative.}} \\
& {\(d = \dfrac{\left|\their{{-\mathbf{i}+2\mathbf{j}-\mathbf{k}}}\right|}{\left|\overrightarrow{AD}\right|} = \dfrac{\sqrt{6}}{\sqrt{3}}\)} & dM1 & {Divides the magnitude of their vector product by the magnitude of their direction vector.} \\
& {\(d = \sqrt{2}\)} & A1 & {Or awrt 1.41.} \\
\nopagebreak & & [4] & \\
8(c) & {\(\overrightarrow{AB}\times\overrightarrow{AD} = \begin{pmatrix}1\\0\\0\end{pmatrix}\times\begin{pmatrix}1\\1\\1\end{pmatrix} = \begin{pmatrix}0\\-1\\1\end{pmatrix}\)} & M1 A1 & {\msNote{M1}{Attempts the vector product of two non-parallel vectors in plane \(ABD\), or finds a normal from two scalar products set to zero. Condone sign slips.}\\ \msNote{A1}{Any nonzero multiple.}} \\
& {\(\mathbf{n}_{ACD} = \begin{pmatrix}-1\\2\\-1\end{pmatrix}\)} & B1 & {Any nonzero multiple; may be taken from (b).} \\
& {\(\cos\theta = \dfrac{\left|\their{(0, -1, 1)}\cdot\their{(-1, 2, -1)}\right|}{\sqrt{2}\,\sqrt{6}}\)} & M1 & {Uses the scalar product of their two normals divided by the product of their magnitudes.} \\
& {\(\cos\theta = \dfrac{3}{\sqrt{12}} = \dfrac{\sqrt{3}}{2}\)\\[4pt] so the acute angle is \(\cos^{-1}\left(\frac{\sqrt{3}}{2}\right)\)} & A1~AG & {cso, from correct normals. A negative cosine must be made positive, or the obtuse angle replaced by its supplement. Needs a concluding statement.} \\
\nopagebreak & & [5] & \\
\nopagebreak \SetCell[c=2]{r} Total for Question 8 & & 16 & \\
\end{msscheme}

\end{document}
